\IfFileExists{stacks-project.cls}{%
\documentclass{stacks-project}
}{%
\documentclass{amsart}
}

% For dealing with references we use the comment environment
\usepackage{verbatim}
\newenvironment{reference}{\comment}{\endcomment}
%\newenvironment{reference}{}{}
\newenvironment{slogan}{\comment}{\endcomment}
\newenvironment{history}{\comment}{\endcomment}

% For commutative diagrams we use Xy-pic
\usepackage[all]{xy}

% We use 2cell for 2-commutative diagrams.
\xyoption{2cell}
\UseAllTwocells

% We use multicol for the list of chapters between chapters
\usepackage{multicol}

% This is generally recommended for better output
\usepackage{lmodern}
\usepackage[T1]{fontenc}

% For cross-file-references
\usepackage{xr-hyper}

% Package for hypertext links:
\usepackage{hyperref}

% For any local file, say "hello.tex" you want to link to please
% use \externaldocument[hello-]{hello}
\externaldocument[introduction-]{introduction}
\externaldocument[conventions-]{conventions}
\externaldocument[sets-]{sets}
\externaldocument[categories-]{categories}
\externaldocument[topology-]{topology}
\externaldocument[sheaves-]{sheaves}
\externaldocument[sites-]{sites}
\externaldocument[stacks-]{stacks}
\externaldocument[fields-]{fields}
\externaldocument[algebra-]{algebra}
\externaldocument[brauer-]{brauer}
\externaldocument[homology-]{homology}
\externaldocument[derived-]{derived}
\externaldocument[simplicial-]{simplicial}
\externaldocument[more-algebra-]{more-algebra}
\externaldocument[smoothing-]{smoothing}
\externaldocument[modules-]{modules}
\externaldocument[sites-modules-]{sites-modules}
\externaldocument[injectives-]{injectives}
\externaldocument[cohomology-]{cohomology}
\externaldocument[sites-cohomology-]{sites-cohomology}
\externaldocument[dga-]{dga}
\externaldocument[dpa-]{dpa}
\externaldocument[sdga-]{sdga}
\externaldocument[hypercovering-]{hypercovering}
\externaldocument[schemes-]{schemes}
\externaldocument[constructions-]{constructions}
\externaldocument[properties-]{properties}
\externaldocument[morphisms-]{morphisms}
\externaldocument[coherent-]{coherent}
\externaldocument[divisors-]{divisors}
\externaldocument[limits-]{limits}
\externaldocument[varieties-]{varieties}
\externaldocument[topologies-]{topologies}
\externaldocument[descent-]{descent}
\externaldocument[perfect-]{perfect}
\externaldocument[more-morphisms-]{more-morphisms}
\externaldocument[flat-]{flat}
\externaldocument[groupoids-]{groupoids}
\externaldocument[more-groupoids-]{more-groupoids}
\externaldocument[etale-]{etale}
\externaldocument[chow-]{chow}
\externaldocument[intersection-]{intersection}
\externaldocument[pic-]{pic}
\externaldocument[weil-]{weil}
\externaldocument[adequate-]{adequate}
\externaldocument[dualizing-]{dualizing}
\externaldocument[duality-]{duality}
\externaldocument[discriminant-]{discriminant}
\externaldocument[derham-]{derham}
\externaldocument[local-cohomology-]{local-cohomology}
\externaldocument[algebraization-]{algebraization}
\externaldocument[curves-]{curves}
\externaldocument[resolve-]{resolve}
\externaldocument[models-]{models}
\externaldocument[functors-]{functors}
\externaldocument[equiv-]{equiv}
\externaldocument[pione-]{pione}
\externaldocument[etale-cohomology-]{etale-cohomology}
\externaldocument[proetale-]{proetale}
\externaldocument[relative-cycles-]{relative-cycles}
\externaldocument[more-etale-]{more-etale}
\externaldocument[trace-]{trace}
\externaldocument[crystalline-]{crystalline}
\externaldocument[spaces-]{spaces}
\externaldocument[spaces-properties-]{spaces-properties}
\externaldocument[spaces-morphisms-]{spaces-morphisms}
\externaldocument[decent-spaces-]{decent-spaces}
\externaldocument[spaces-cohomology-]{spaces-cohomology}
\externaldocument[spaces-limits-]{spaces-limits}
\externaldocument[spaces-divisors-]{spaces-divisors}
\externaldocument[spaces-over-fields-]{spaces-over-fields}
\externaldocument[spaces-topologies-]{spaces-topologies}
\externaldocument[spaces-descent-]{spaces-descent}
\externaldocument[spaces-perfect-]{spaces-perfect}
\externaldocument[spaces-more-morphisms-]{spaces-more-morphisms}
\externaldocument[spaces-flat-]{spaces-flat}
\externaldocument[spaces-groupoids-]{spaces-groupoids}
\externaldocument[spaces-more-groupoids-]{spaces-more-groupoids}
\externaldocument[bootstrap-]{bootstrap}
\externaldocument[spaces-pushouts-]{spaces-pushouts}
\externaldocument[spaces-chow-]{spaces-chow}
\externaldocument[groupoids-quotients-]{groupoids-quotients}
\externaldocument[spaces-more-cohomology-]{spaces-more-cohomology}
\externaldocument[spaces-simplicial-]{spaces-simplicial}
\externaldocument[spaces-duality-]{spaces-duality}
\externaldocument[formal-spaces-]{formal-spaces}
\externaldocument[restricted-]{restricted}
\externaldocument[spaces-resolve-]{spaces-resolve}
\externaldocument[formal-defos-]{formal-defos}
\externaldocument[defos-]{defos}
\externaldocument[cotangent-]{cotangent}
\externaldocument[examples-defos-]{examples-defos}
\externaldocument[algebraic-]{algebraic}
\externaldocument[examples-stacks-]{examples-stacks}
\externaldocument[stacks-sheaves-]{stacks-sheaves}
\externaldocument[criteria-]{criteria}
\externaldocument[artin-]{artin}
\externaldocument[quot-]{quot}
\externaldocument[stacks-properties-]{stacks-properties}
\externaldocument[stacks-morphisms-]{stacks-morphisms}
\externaldocument[stacks-limits-]{stacks-limits}
\externaldocument[stacks-cohomology-]{stacks-cohomology}
\externaldocument[stacks-perfect-]{stacks-perfect}
\externaldocument[stacks-introduction-]{stacks-introduction}
\externaldocument[stacks-more-morphisms-]{stacks-more-morphisms}
\externaldocument[stacks-geometry-]{stacks-geometry}
\externaldocument[moduli-]{moduli}
\externaldocument[moduli-curves-]{moduli-curves}
\externaldocument[examples-]{examples}
\externaldocument[exercises-]{exercises}
\externaldocument[guide-]{guide}
\externaldocument[desirables-]{desirables}
\externaldocument[coding-]{coding}
\externaldocument[obsolete-]{obsolete}
\externaldocument[fdl-]{fdl}
\externaldocument[index-]{index}

% Theorem environments.
%
\theoremstyle{plain}
\newtheorem{theorem}[subsection]{Theorem}
\newtheorem{proposition}[subsection]{Proposition}
\newtheorem{lemma}[subsection]{Lemma}

\theoremstyle{definition}
\newtheorem{definition}[subsection]{Definition}
\newtheorem{example}[subsection]{Example}
\newtheorem{exercise}[subsection]{Exercise}
\newtheorem{situation}[subsection]{Situation}

\theoremstyle{remark}
\newtheorem{remark}[subsection]{Remark}
\newtheorem{remarks}[subsection]{Remarks}

\numberwithin{equation}{subsection}

% Macros
%
\def\lim{\mathop{\mathrm{lim}}\nolimits}
\def\colim{\mathop{\mathrm{colim}}\nolimits}
\def\Spec{\mathop{\mathrm{Spec}}}
\def\Hom{\mathop{\mathrm{Hom}}\nolimits}
\def\Ext{\mathop{\mathrm{Ext}}\nolimits}
\def\SheafHom{\mathop{\mathcal{H}\!\mathit{om}}\nolimits}
\def\SheafExt{\mathop{\mathcal{E}\!\mathit{xt}}\nolimits}
\def\Sch{\mathit{Sch}}
\def\Mor{\mathop{\mathrm{Mor}}\nolimits}
\def\Ob{\mathop{\mathrm{Ob}}\nolimits}
\def\Sh{\mathop{\mathit{Sh}}\nolimits}
\def\NL{\mathop{N\!L}\nolimits}
\def\CH{\mathop{\mathrm{CH}}\nolimits}
\def\proetale{{pro\text{-}\acute{e}tale}}
\def\etale{{\acute{e}tale}}
\def\QCoh{\mathit{QCoh}}
\def\Ker{\mathop{\mathrm{Ker}}}
\def\Im{\mathop{\mathrm{Im}}}
\def\Coker{\mathop{\mathrm{Coker}}}
\def\Coim{\mathop{\mathrm{Coim}}}

% Boxtimes
%
\DeclareMathSymbol{\boxtimes}{\mathbin}{AMSa}{"02}

%
% Macros for moduli stacks/spaces
%
\def\QCohstack{\mathcal{QC}\!\mathit{oh}}
\def\Cohstack{\mathcal{C}\!\mathit{oh}}
\def\Spacesstack{\mathcal{S}\!\mathit{paces}}
\def\Quotfunctor{\mathrm{Quot}}
\def\Hilbfunctor{\mathrm{Hilb}}
\def\Curvesstack{\mathcal{C}\!\mathit{urves}}
\def\Polarizedstack{\mathcal{P}\!\mathit{olarized}}
\def\Complexesstack{\mathcal{C}\!\mathit{omplexes}}
% \Pic is the operator that assigns to X its picard group, usage \Pic(X)
% \Picardstack_{X/B} denotes the Picard stack of X over B
% \Picardfunctor_{X/B} denotes the Picard functor of X over B
\def\Pic{\mathop{\mathrm{Pic}}\nolimits}
\def\Picardstack{\mathcal{P}\!\mathit{ic}}
\def\Picardfunctor{\mathrm{Pic}}
\def\Deformationcategory{\mathcal{D}\!\mathit{ef}}
\usepackage{cleveref}
\crefname{originallemma}{original lemma}{original lemmas}
\Crefname{originallemma}{Original lemma}{Original lemmas}
\makeatletter
\newlabel{morphisms-lemma-finite-fibre}{{02NG}{source}{}{sourcepointer-1}{}}
\newlabel{morphisms-lemma-finite-fibre@cref}{{[originallemma][0][]02NG}{[1][0][]source}{}{}{}}
\makeatother
\crefname{originallemma}{original lemma}{original lemmas}
\Crefname{originallemma}{Original lemma}{Original lemmas}
\makeatletter
\newlabel{morphisms-lemma-composition-quasi-finite}{{01TL}{source}{}{sourcepointer-2}{}}
\newlabel{morphisms-lemma-composition-quasi-finite@cref}{{[originallemma][0][]01TL}{[1][0][]source}{}{}{}}
\makeatother
\crefname{originallemma}{original lemma}{original lemmas}
\Crefname{originallemma}{Original lemma}{Original lemmas}
\makeatletter
\newlabel{morphisms-lemma-quasi-finite-at-point-characterize}{{01TH}{source}{}{sourcepointer-3}{}}
\newlabel{morphisms-lemma-quasi-finite-at-point-characterize@cref}{{[originallemma][0][]01TH}{[1][0][]source}{}{}{}}
\makeatother
\crefname{originallemma}{original lemma}{original lemmas}
\Crefname{originallemma}{Original lemma}{Original lemmas}
\makeatletter
\newlabel{morphisms-lemma-base-change-flat}{{01U9}{source}{}{sourcepointer-4}{}}
\newlabel{morphisms-lemma-base-change-flat@cref}{{[originallemma][0][]01U9}{[1][0][]source}{}{}{}}
\makeatother
\crefname{originallemma}{original lemma}{original lemmas}
\Crefname{originallemma}{Original lemma}{Original lemmas}
\makeatletter
\newlabel{morphisms-lemma-base-change-finite-presentation}{{01TS}{source}{}{sourcepointer-5}{}}
\newlabel{morphisms-lemma-base-change-finite-presentation@cref}{{[originallemma][0][]01TS}{[1][0][]source}{}{}{}}
\makeatother
\crefname{originallemma}{original lemma}{original lemmas}
\Crefname{originallemma}{Original lemma}{Original lemmas}
\makeatletter
\newlabel{morphisms-lemma-fppf-open}{{01UA}{source}{}{sourcepointer-6}{}}
\newlabel{morphisms-lemma-fppf-open@cref}{{[originallemma][0][]01UA}{[1][0][]source}{}{}{}}
\makeatother
\crefname{originallemma}{original lemma}{original lemmas}
\Crefname{originallemma}{Original lemma}{Original lemmas}
\makeatletter
\newlabel{morphisms-lemma-locally-quasi-finite-qc-source-universally-bounded}{{03JA}{source}{}{sourcepointer-7}{}}
\newlabel{morphisms-lemma-locally-quasi-finite-qc-source-universally-bounded@cref}{{[originallemma][0][]03JA}{[1][0][]source}{}{}{}}
\makeatother
\crefname{originallemma}{original lemma}{original lemmas}
\Crefname{originallemma}{Original lemma}{Original lemmas}
\makeatletter
\newlabel{lemma-two-fibres}{{lemma-two-fibres}{source}{}{sourcepointer-8}{}}
\newlabel{lemma-two-fibres@cref}{{[originallemma][0][]lemma-two-fibres}{[1][0][]source}{}{}{}}
\makeatother
\crefname{originallemma}{original lemma}{original lemmas}
\Crefname{originallemma}{Original lemma}{Original lemmas}
\makeatletter
\newlabel{more-morphisms-lemma-etale-splits-off-quasi-finite-part-technical}{{02LN}{source}{}{sourcepointer-9}{}}
\newlabel{more-morphisms-lemma-etale-splits-off-quasi-finite-part-technical@cref}{{[originallemma][0][]02LN}{[1][0][]source}{}{}{}}
\makeatother
\title{Etale neighbourhood splitting a finite groupoid piece off a quasi-finite point}
\author{The Stacks Project Authors}
\date{stacks-project:more-groupoids.tex:local-snapshot-20261012; GNU FDL 1.2 or later}
\begin{document}
\maketitle
\setcounter{subsection}{1}
\noindent\textit{Editorial note.} All delivered mathematical text is the author's own, copied from the pinned Stacks Project chapter file more-groupoids.tex; no statement or proof step was rewritten, reordered or removed. The counted span is the author environment and its own proof environment. References inside the span are handed to the assembler as navigation stubs that carry the canonical tag of the cited tag only; whether each of them is load-bearing is recorded on the stub itself.\par\medskip
\noindent\textit{Allowed standard prerequisites (editorial).}\begin{itemize}
\item Quasi-finiteness at the points of a finite fibre (morphisms chapter, Lemma morphisms-lemma-finite-fibre, canonical tag 02NG). cited at proof lines 2450-2452. The structure map s is quasi-finite at each of the finitely many points of the source fibre of u.
\item Quasi-finiteness is preserved under composition (morphisms chapter, Lemma morphisms-lemma-composition-quasi-finite, canonical tag 01TL). cited at proof lines 2453-2454. The composite structure morphism from R to S is quasi-finite at each r\_i.
\item Characterization of quasi-finite at a point by isolation in the fibre (morphisms chapter, Lemma morphisms-lemma-quasi-finite-at-point-characterize, canonical tag 01TH). cited at proof lines 2455-2456 and 2468-2470. Isolates each r\_i in the fibre of R over p and, together with openness of the target map of fibres, isolates each u\_j in the fibre of U over p, giving quasi-finiteness of the structure morphism at the u\_j.
\item Flatness is preserved under base change (morphisms chapter, Lemma morphisms-lemma-base-change-flat, canonical tag 01U9). cited at proof lines 2460-2464. The target map of the fibres over p is flat.
\item Being locally of finite presentation is preserved under base change (morphisms chapter, Lemma morphisms-lemma-base-change-finite-presentation, canonical tag 01TS). cited at proof lines 2460-2464. The target map of the fibres over p is locally of finite presentation.
\item Flat morphisms locally of finite presentation are open (morphisms chapter, Lemma morphisms-lemma-fppf-open, canonical tag 01UA). cited at proof lines 2465-2466. The target map of fibres is open, which propagates isolation from the r\_i to the points u\_j = t(r\_i).
\item Finiteness of scheme-theoretic fibres with finitely many points for locally finite type morphisms, much more general form (morphisms chapter, Lemma morphisms-lemma-locally-quasi-finite-qc-source-universally-bounded, canonical tag 03JA). cited at proof lines 2474-2478. Supports that the scheme-theoretic source fibre F\_u is finite over the residue field of u.
\item Scheme-theoretic source fibres over two points joined by an arrow with common target become isomorphic over a common field extension (more-groupoids chapter, Lemma lemma-two-fibres, no canonical tag in the snapshot; no-tag pointer to the source label line). cited at proof lines 2480-2482. Transfers finiteness of F\_u over the residue field of u to finiteness of F\_\{u\_j\} over the residue field of u\_j, extending the quasi-finite setup from u to every object of its orbit.
\item Technical etale splitting of a quasi-finite part: over an etale neighbourhood with residue field identification the base change decomposes into an open and closed finite piece whose fibre maps bijectively onto the prescribed points and an open complement avoiding them (more-morphisms chapter, Lemma more-morphisms-lemma-etale-splits-off-quasi-finite-part-technical, canonical tag 02LN). cited at proof lines 2506-2521. Produces the etale neighbourhood and the decompositions of the base changes into U' disjoint union W and R' disjoint union W' with U' and R' finite over S', whence conclusions (a) through (d) and the fibre level versions of (e) and (f).
\end{itemize}\medskip
\ifdefined\hypertarget\hypertarget{sourcepointer-1}{}\fi
\noindent\textit{Source reference.} Original lemma 02NG: sources/stacks-chow/tags.txt, source lines 3479--3479. Cited at proof lines 2450-2452: s is quasi-finite at each point of the finite source fibre of u.\par\medskip
\ifdefined\hypertarget\hypertarget{sourcepointer-2}{}\fi
\noindent\textit{Source reference.} Original lemma 01TL: sources/stacks-chow/tags.txt, source lines 2378--2378. Cited at proof lines 2453-2454: the composite to S is quasi-finite at each r\_i.\par\medskip
\ifdefined\hypertarget\hypertarget{sourcepointer-3}{}\fi
\noindent\textit{Source reference.} Original lemma 01TH: sources/stacks-chow/tags.txt, source lines 2374--2374. Cited at proof lines 2455-2456 and 2468-2470: isolates r\_i in the fibre of R over p and the u\_j in the fibre of U over p.\par\medskip
\ifdefined\hypertarget\hypertarget{sourcepointer-4}{}\fi
\noindent\textit{Source reference.} Original lemma 01U9: sources/stacks-chow/tags.txt, source lines 2402--2402. Cited at proof lines 2460-2464: the target map of the fibres over p is flat.\par\medskip
\ifdefined\hypertarget\hypertarget{sourcepointer-5}{}\fi
\noindent\textit{Source reference.} Original lemma 01TS: sources/stacks-chow/tags.txt, source lines 2385--2385. Cited at proof lines 2460-2464: the target map of the fibres over p is locally of finite presentation.\par\medskip
\ifdefined\hypertarget\hypertarget{sourcepointer-6}{}\fi
\noindent\textit{Source reference.} Original lemma 01UA: sources/stacks-chow/tags.txt, source lines 2403--2403. Cited at proof lines 2465-2466: the target map of fibres is open, propagating isolation from r\_i to u\_j.\par\medskip
\ifdefined\hypertarget\hypertarget{sourcepointer-7}{}\fi
\noindent\textit{Source reference.} Original lemma 03JA: sources/stacks-chow/tags.txt, source lines 4632--4632. Cited at proof lines 2474-2478: supports finiteness of the scheme-theoretic fibre F\_u over the residue field of u.\par\medskip
\ifdefined\hypertarget\hypertarget{sourcepointer-8}{}\fi
\noindent\textit{Source reference.} Original lemma lemma-two-fibres: sources/stacks-more-groupoids/pinned-original-source.tex, source lines 761--761. Cited at proof lines 2480-2482: identifies F\_u and F\_\{u\_j\} over a common field extension, transferring fibre finiteness within the orbit.\par\medskip
\ifdefined\hypertarget\hypertarget{sourcepointer-9}{}\fi
\noindent\textit{Source reference.} Original lemma 02LN: sources/stacks-chow/tags.txt, source lines 3414--3414. Cited at proof lines 2506-2521: produces the etale neighbourhood and the decompositions into finite pieces with bijective fibres.\par\medskip
\begin{lemma}
\label{lemma-quasi-finite-over-base}
Let $S$ be a scheme.
Let $(U, R, s, t, c)$ be a groupoid scheme over $S$.
Let $p \in S$ be a point, and let $u \in U$ be a point lying over $p$.
Assume that
\begin{enumerate}
\item $U \to S$ is locally of finite type,
\item $U \to S$ is quasi-finite at $u$,
\item $U \to S$ is separated,
\item $R \to S$ is separated,
\item $s$, $t$ are flat and locally of finite presentation, and
\item $s^{-1}(\{u\})$ is finite.
\end{enumerate}
Then there exists an \'etale neighbourhood $(S', p') \to (S, p)$ with
$\kappa(p) = \kappa(p')$ and a base change diagram
$$
\xymatrix{
R' \amalg W'
\ar@{=}[r] &
S' \times_S R
\ar[r] \ar@<2ex>[d]^{s'} \ar@<-2ex>[d]_{t'} &
R \ar@<1ex>[d]^s \ar@<-1ex>[d]_t \\
U' \amalg W
\ar@{=}[r] &
S' \times_S U
\ar[r] \ar[d] &
U \ar[d] \\
 &
S' \ar[r] &
S
}
$$
where the equal signs are decompositions into open and closed
subschemes such that
\begin{enumerate}
\item[(a)] there exists a point $u'$ of $U'$ mapping to $u$ in $U$,
\item[(b)] the fibre $(U')_{p'}$ equals $t'\big((s')^{-1}(\{u'\})\big)$
set theoretically,
\item[(c)] the fibre $(R')_{p'}$ equals $(s')^{-1}\big((U')_{p'}\big)$
set theoretically,
\item[(d)] the schemes $U'$ and $R'$ are finite over $S'$,
\item[(e)] we have $s'(R') \subset U'$ and $t'(R') \subset U'$,
\item[(f)] we have
$c'(R' \times_{s', U', t'} R') \subset R'$
where $c'$ is the base change of $c$, and
\item[(g)] the morphisms $s', t', c'$ determine a groupoid structure
by taking the system
$(U', R', s'|_{R'}, t'|_{R'}, c'|_{R' \times_{s', U', t'} R'})$.
\end{enumerate}
\end{lemma}

\begin{proof}
Let us denote $f : U \to S$ the structure morphism of $U$.
By assumption (6) we can write $s^{-1}(\{u\}) = \{r_1, \ldots, r_n\}$.
Since this set is finite, we see that $s$ is quasi-finite at each of
these finitely many inverse images, see
Morphisms, Lemma \ref{morphisms-lemma-finite-fibre}.
Hence we see that $f \circ s : R \to S$ is quasi-finite at each $r_i$
(Morphisms, Lemma \ref{morphisms-lemma-composition-quasi-finite}).
Hence $r_i$ is isolated in the fibre $R_p$, see
Morphisms, Lemma \ref{morphisms-lemma-quasi-finite-at-point-characterize}.
Write $t(\{r_1, \ldots, r_n\}) = \{u_1, \ldots, u_m\}$.
Note that it may happen that $m < n$ and note that
$u \in \{u_1, \ldots, u_m\}$.
Since $t$ is flat and locally of finite presentation,
the morphism of fibres $t_p : R_p \to U_p$ is flat and locally of
finite presentation (Morphisms,
Lemmas \ref{morphisms-lemma-base-change-flat} and
\ref{morphisms-lemma-base-change-finite-presentation}),
hence open (Morphisms,
Lemma \ref{morphisms-lemma-fppf-open}).
The fact that each $r_i$ is isolated in $R_p$ implies that
each $u_j = t(r_i)$ is isolated in $U_p$. Using
Morphisms, Lemma \ref{morphisms-lemma-quasi-finite-at-point-characterize}
again, we see that $f$ is quasi-finite at $u_1, \ldots, u_m$.

\medskip\noindent
Denote $F_u = s^{-1}(u)$ and $F_{u_j} = s^{-1}(u_j)$ the scheme theoretic
fibres. Note that $F_u$ is finite over $\kappa(u)$ as it is locally of finite
type over $\kappa(u)$ with finitely many points (for example it follows from
the much more general
Morphisms, Lemma
\ref{morphisms-lemma-locally-quasi-finite-qc-source-universally-bounded}).
By
Lemma \ref{lemma-two-fibres}
we see that $F_u$ and $F_{u_j}$ become isomorphic over a common
field extension of $\kappa(u)$ and $\kappa(u_j)$. Hence we see
that $F_{u_j}$ is finite over $\kappa(u_j)$. In particular we see
$s^{-1}(\{u_j\})$ is a finite set for each $j = 1, \ldots, m$.
Thus we see that assumptions (2) and (6) hold for each $u_j$ also
(above we saw that $U \to S$ is quasi-finite at $u_j$).
Hence the argument of the first paragraph applies to each $u_j$
and we see that $R \to U$ is quasi-finite at each of the points of
$$
\{r_1, \ldots, r_N\} = s^{-1}(\{u_1, \ldots, u_m\})
$$
Note that $t(\{r_1, \ldots, r_N\}) = \{u_1, \ldots, u_m\}$ and
$t^{-1}(\{u_1, \ldots, u_m\}) = \{r_1, \ldots, r_N\}$
since $R$ is a groupoid\footnote{Explanation in groupoid language:
The original set $\{r_1, \ldots, r_n\}$ was the set of arrows with
source $u$. The set $\{u_1, \ldots, u_m\}$ was the set of objects
isomorphic to $u$. And $\{r_1, \ldots, r_N\}$ is the set of all arrows
between all the objects equivalent to $u$.}. Moreover, we have
$\text{pr}_0(c^{-1}(\{r_1, \ldots, r_N\})) = \{r_1, \ldots, r_N\}$
and
$\text{pr}_1(c^{-1}(\{r_1, \ldots, r_N\})) = \{r_1, \ldots, r_N\}$.
Similarly we get $e(\{u_1, \ldots, u_m\}) \subset \{r_1, \ldots, r_N\}$
and $i(\{r_1, \ldots, r_N\}) = \{r_1, \ldots, r_N\}$.

\medskip\noindent
We may apply
More on Morphisms,
Lemma \ref{more-morphisms-lemma-etale-splits-off-quasi-finite-part-technical}
to the pairs
$(U \to S, \{u_1, \ldots, u_m\})$ and
$(R \to S, \{r_1, \ldots, r_N\})$
to get an \'etale neighbourhood $(S', p') \to (S, p)$
which induces an identification $\kappa(p) = \kappa(p')$
such that $S' \times_S U$ and $S' \times_S R$ decompose as
$$
S' \times_S U = U' \amalg W, \quad
S' \times_S R = R' \amalg W'
$$
with $U' \to S'$ finite and $(U')_{p'}$ mapping bijectively to
$\{u_1, \ldots, u_m\}$, and $R' \to S'$ finite and
$(R')_{p'}$ mapping bijectively to $\{r_1, \ldots, r_N\}$.
Moreover, no point of $W_{p'}$ (resp.\ $(W')_{p'}$) maps to
any of the points $u_j$ (resp.\ $r_i$). At this point (a), (b), (c), and (d)
of the lemma are satisfied. Moreover, the inclusions of (e) and (f) hold
on fibres over $p'$, i.e., $s'((R')_{p'}) \subset (U')_{p'}$,
$t'((R')_{p'}) \subset (U')_{p'}$, and
$c'((R' \times_{s', U', t'} R')_{p'}) \subset (R')_{p'}$.

\medskip\noindent
We claim that we can replace $S'$ by a Zariski open neighbourhood
of $p'$ so that the inclusions of (e) and (f) hold.
For example, consider the set $E = (s'|_{R'})^{-1}(W)$.
This is open and closed in $R'$ and does not contain any points
of $R'$ lying over $p'$. Since $R' \to S'$ is closed,
after replacing $S'$ by $S' \setminus (R' \to S')(E)$ we reach a
situation where $E$ is empty. In other words $s'$ maps $R'$ into $U'$.
Note that this property is preserved under further shrinking $S'$.
Similarly, we can arrange it so that $t'$ maps $R'$ into $U'$.
At this point (e) holds. In the same manner, consider the set
$E = (c'|_{R' \times_{s', U', t'} R'})^{-1}(W')$.
It is open and closed in the scheme $R' \times_{s', U', t'} R'$
which is finite over $S'$, and does not contain any points lying
over $p'$. Hence after replacing $S'$ by
$S' \setminus (R' \times_{s', U', t'} R' \to S')(E)$
we reach a situation where $E$ is empty. In other words we obtain
the inclusion in (f). We may repeat the argument also with the identity
$e' : S' \times_S U \to S' \times_S R$ and the inverse
$i' : S' \times_S R \to S' \times_S R$ so that we may assume
(after shrinking $S'$ some more) that $(e'|_{U'})^{-1}(W') = \emptyset$
and $(i'|_{R'})^{-1}(W') = \emptyset$.

\medskip\noindent
At this point we see that we may consider the structure
$$
(U', R', s'|_{R'}, t'|_{R'}, c'|_{R' \times_{t', U', s'} R'},
e'|_{U'}, i'|_{R'}).
$$
The axioms of a groupoid scheme over $S'$ hold
because they hold for the groupoid scheme
$(S' \times_S U, S' \times_S R, s', t', c', e', i')$.
\end{proof}

\end{document}
